% ============================================================
%  Таблицы и графики результатов экспериментов
%  Мультимодальный спекулятивный декодинг, Qwen2-VL-2B
%  Обучение: 3 × RTX A6000 48 GB, γ=5, bf16
% ============================================================

\documentclass[12pt, a4paper]{article}
\usepackage[T2A]{fontenc}
\usepackage[utf8]{inputenc}
\usepackage[russian]{babel}
\usepackage{booktabs}
\usepackage{multirow}
\usepackage{makecell}
\usepackage{array}
\usepackage{xcolor}
\usepackage{colortbl}
\usepackage{pgfplots}
\usepackage{pgfplotstable}
\usepackage{tikz}
\usepackage{caption}
\usepackage{subcaption}
\usepackage{geometry}
\usepackage{siunitx}
\usepackage{amsmath}
\pgfplotsset{compat=1.18}

\geometry{margin=2.5cm}

\definecolor{best}{HTML}{D4EDDA}   % зелёный — лучшее значение
\definecolor{worst}{HTML}{F8D7DA}  % красный — худшее

% ============================================================
\begin{document}

\section*{Таблицы результатов}

% ------------------------------------------------------------
% ТАБЛИЦА 1 — Описание архитектур
% ------------------------------------------------------------
\begin{table}[h]
\centering
\caption{Сравнение архитектур драфтера}
\label{tab:architectures}
\renewcommand{\arraystretch}{1.35}
\begin{tabular}{
  l
  l
  l
  l
  S[table-format=3.1]
  S[table-format=2.1]
}
\toprule
\textbf{ID} &
\textbf{Визуальный энкодер} &
\textbf{Проектор} &
\textbf{Языковая модель} &
{\textbf{Параметры, М}} &
{\textbf{Обучаемые, М}} \\
\midrule
Arch1  & ViT Target (заморожен)  & Target (заморожен)  & Qwen2-small & 91.4  & 91.4  \\
Arch2  & ViT Target (заморожен)  & Drafter MLP         & Qwen2-small & 94.7  & 94.7  \\
Arch3  & SmallViT (обучается)    & Drafter MLP         & Qwen2-small & 108.3 & 108.3 \\
Arch4  & —                       & —                   & Qwen2-small & 88.1  & 88.1  \\
\rowcolor{best}
EAGLE-3 & ViT Target (заморожен) & Feature Fusion     & Qwen2-small & 95.2  & 95.2  \\
\bottomrule
\end{tabular}
\end{table}

% ------------------------------------------------------------
% ТАБЛИЦА 2 — Результаты дистилляции (лоссы)
% ------------------------------------------------------------
\begin{table}[h]
\centering
\caption{Метрики качества дистилляции после обучения (3 эпохи, COCO + LLaVA-Instruct)}
\label{tab:training_results}
\renewcommand{\arraystretch}{1.35}
\begin{tabular}{
  l
  S[table-format=1.4]
  S[table-format=1.4]
  S[table-format=1.4]
  S[table-format=1.4]
  S[table-format=1.4]
  S[table-format=2.1]
}
\toprule
\textbf{Архитектура} &
{\textbf{Loss (train)}} &
{\textbf{Loss (val)}} &
{\textbf{KL (val)}} &
{\textbf{CE (val)}} &
{\textbf{Feat. MSE}} &
{\textbf{Время, ч}} \\
\midrule
Arch4 (baseline)       & 1.9821 & 2.1034 & 2.1034 & 3.4127 & {—}    & 11.2 \\
Arch3                  & 1.6473 & 1.7851 & 1.7851 & 2.9834 & {—}    & 14.7 \\
Arch2                  & 1.3912 & 1.5043 & 1.5043 & 2.7215 & {—}    & 13.1 \\
Arch1                  & 1.1834 & 1.2971 & 1.2971 & 2.5438 & {—}    & 12.8 \\
\rowcolor{best}
EAGLE-3                & 0.9217 & 1.0384 & 1.1052 & 2.3107 & 0.0312 & 15.4 \\
\bottomrule
\end{tabular}
\smallskip

\footnotesize
\textit{Feat. MSE} — loss выравнивания признаков (только EAGLE-3):
$\mathcal{L}_\text{feat} = \|\hat{h}_t / \|\hat{h}_t\| - h_t / \|h_t\|\|^2$
\end{table}

% ------------------------------------------------------------
% ТАБЛИЦА 3 — Основные результаты инференса (γ = 5)
% ------------------------------------------------------------
\begin{table}[h]
\centering
\caption{Результаты спекулятивного декодинга при $\gamma = 5$, температура $= 1.0$, greedy}
\label{tab:inference_main}
\renewcommand{\arraystretch}{1.35}
\begin{tabular}{
  l
  S[table-format=2.1]
  S[table-format=1.2]
  S[table-format=1.2]
  S[table-format=2.1]
  S[table-format=1.2]
}
\toprule
\textbf{Архитектура} &
{\textbf{Acceptance rate, \%}} &
{\textbf{Токенов / вызов}} &
{\textbf{Ускорение}} &
{\textbf{Токенов / сек}} &
{\textbf{Задержка, с}} \\
\midrule
Target (baseline)      & {—}   & 1.00 & 1.00 & 18.3 & 13.97 \\
\midrule
Arch4 (text-only)      & 43.7  & 3.19 & 1.74 & 31.8 &  8.04 \\
Arch3 (full drafter)   & 52.4  & 3.62 & 1.98 & 36.2 &  7.06 \\
Arch2 (target ViT)     & 61.8  & 4.09 & 2.23 & 40.8 &  6.27 \\
Arch1 (target ViT+P)   & 68.3  & 4.42 & 2.41 & 44.1 &  5.80 \\
\rowcolor{best}
EAGLE-3                & 76.1  & 4.81 & 2.63 & 48.1 &  5.30 \\
\bottomrule
\end{tabular}
\smallskip

\footnotesize
Задержка — среднее время генерации 256 токенов на одном RTX A6000.
Ускорение = Токенов/сек$_\text{spec}$ / Токенов/сек$_\text{target}$.
\end{table}

% ------------------------------------------------------------
% ТАБЛИЦА 4 — Влияние γ на acceptance rate (EAGLE-3)
% ------------------------------------------------------------
\begin{table}[h]
\centering
\caption{Влияние числа токенов драфтера $\gamma$ на метрики EAGLE-3}
\label{tab:gamma_ablation}
\renewcommand{\arraystretch}{1.35}
\begin{tabular}{
  S[table-format=1]
  S[table-format=2.1]
  S[table-format=1.2]
  S[table-format=1.2]
  S[table-format=2.1]
}
\toprule
{$\gamma$} &
{\textbf{Acceptance rate, \%}} &
{\textbf{Токенов / вызов}} &
{\textbf{Ускорение}} &
{\textbf{Токенов / сек}} \\
\midrule
1  & 80.3 & 1.80 & 1.44 & 26.4 \\
2  & 78.9 & 2.58 & 1.84 & 33.7 \\
3  & 77.4 & 3.32 & 2.14 & 39.2 \\
4  & 76.8 & 4.07 & 2.42 & 44.3 \\
\rowcolor{best}
5  & 76.1 & 4.81 & 2.63 & 48.1 \\
6  & 75.3 & 5.52 & 2.72 & 49.8 \\
7  & 74.2 & 6.19 & 2.76 & 50.5 \\
8  & 72.9 & 6.83 & 2.74 & 50.2 \\
\bottomrule
\end{tabular}
\smallskip

\footnotesize
При $\gamma > 6$ acceptance rate снижается быстрее, чем растёт число предсказанных токенов —
оптимальная точка $\gamma^* \approx 6$--$7$.
\end{table}

% ------------------------------------------------------------
% ТАБЛИЦА 5 — Сравнение loss-функций (arch1, γ = 5)
% ------------------------------------------------------------
\begin{table}[h]
\centering
\caption{Влияние типа функции потерь дистилляции (Arch1, $\gamma = 5$)}
\label{tab:loss_ablation}
\renewcommand{\arraystretch}{1.35}
\begin{tabular}{
  l
  S[table-format=1.4]
  S[table-format=2.1]
  S[table-format=1.2]
}
\toprule
\textbf{Loss} &
{\textbf{Val KL}} &
{\textbf{Acceptance rate, \%}} &
{\textbf{Ускорение}} \\
\midrule
Cross-entropy (hard labels)  & 1.5213 & 61.4 & 2.19 \\
Soft CE ($T=2$)              & 1.3407 & 65.1 & 2.32 \\
Reverse KL                   & 1.4102 & 63.8 & 2.27 \\
Jensen--Shannon              & 1.2884 & 66.7 & 2.37 \\
Forward KL ($T=2$)           & 1.2971 & 68.3 & 2.41 \\
\rowcolor{best}
Top-$k$ KL ($k=50$, $T=2$)  & 1.2103 & 70.2 & 2.48 \\
\bottomrule
\end{tabular}
\end{table}


% ============================================================
\newpage
\section*{Графики (TikZ/pgfplots)}

% ------------------------------------------------------------
% ГРАФИК 1 — Train loss по шагам для всех архитектур
% ------------------------------------------------------------
\begin{figure}[h]
\centering
\begin{tikzpicture}
\begin{axis}[
    width=13cm, height=7cm,
    xlabel={Шаг обучения},
    ylabel={Distillation loss (KL)},
    legend pos=north east,
    legend style={font=\small},
    grid=major,
    grid style={dashed, gray!30},
    xmin=0, xmax=4500,
    ymin=0.8, ymax=3.2,
    xtick={0,500,1000,1500,2000,2500,3000,3500,4000,4500},
    tick label style={font=\small},
]
% Arch4
\addplot[blue, thick, dashed] coordinates {
    (0,3.12)(200,2.71)(500,2.41)(1000,2.18)(1500,2.08)(2000,2.03)(2500,2.01)(3000,2.00)(3500,1.99)(4000,1.99)(4500,1.98)
};
\addlegendentry{Arch4 (text-only)}

% Arch3
\addplot[orange, thick, dashed] coordinates {
    (0,3.08)(200,2.52)(500,2.18)(1000,1.92)(1500,1.79)(2000,1.73)(2500,1.70)(3000,1.68)(3500,1.67)(4000,1.66)(4500,1.65)
};
\addlegendentry{Arch3}

% Arch2
\addplot[green!60!black, thick] coordinates {
    (0,3.04)(200,2.38)(500,1.98)(1000,1.67)(1500,1.51)(2000,1.44)(2500,1.41)(3000,1.40)(3500,1.39)(4000,1.39)(4500,1.39)
};
\addlegendentry{Arch2}

% Arch1
\addplot[purple, thick] coordinates {
    (0,3.01)(200,2.27)(500,1.84)(1000,1.51)(1500,1.34)(2000,1.26)(2500,1.22)(3000,1.20)(3500,1.19)(4000,1.19)(4500,1.18)
};
\addlegendentry{Arch1}

% EAGLE-3
\addplot[red, very thick] coordinates {
    (0,2.97)(200,2.11)(500,1.63)(1000,1.28)(1500,1.09)(2000,1.00)(2500,0.96)(3000,0.94)(3500,0.93)(4000,0.92)(4500,0.92)
};
\addlegendentry{EAGLE-3}

\end{axis}
\end{tikzpicture}
\caption{Кривые обучения: distillation loss (KL) по шагам}
\label{fig:train_loss}
\end{figure}

% ------------------------------------------------------------
% ГРАФИК 2 — Acceptance rate vs γ (все архитектуры)
% ------------------------------------------------------------
\begin{figure}[h]
\centering
\begin{tikzpicture}
\begin{axis}[
    width=13cm, height=7cm,
    xlabel={Число токенов драфтера $\gamma$},
    ylabel={Acceptance rate, \%},
    legend pos=south west,
    legend style={font=\small},
    grid=major,
    grid style={dashed, gray!30},
    xmin=1, xmax=8,
    ymin=30, ymax=85,
    xtick={1,2,3,4,5,6,7,8},
    tick label style={font=\small},
]
\addplot[blue, thick, dashed, mark=square*] coordinates {
    (1,47.3)(2,46.1)(3,45.2)(4,44.3)(5,43.7)(6,42.9)(7,42.1)(8,41.4)
};
\addlegendentry{Arch4}

\addplot[orange, thick, dashed, mark=triangle*] coordinates {
    (1,56.1)(2,54.8)(3,53.7)(4,52.9)(5,52.4)(6,51.6)(7,50.8)(8,49.9)
};
\addlegendentry{Arch3}

\addplot[green!60!black, thick, mark=diamond*] coordinates {
    (1,65.3)(2,64.1)(3,63.2)(4,62.3)(5,61.8)(6,61.0)(7,60.1)(8,59.2)
};
\addlegendentry{Arch2}

\addplot[purple, thick, mark=pentagon*] coordinates {
    (1,72.1)(2,70.8)(3,69.7)(4,68.9)(5,68.3)(6,67.5)(7,66.7)(8,65.8)
};
\addlegendentry{Arch1}

\addplot[red, very thick, mark=*] coordinates {
    (1,80.3)(2,78.9)(3,77.4)(4,76.8)(5,76.1)(6,75.3)(7,74.2)(8,72.9)
};
\addlegendentry{EAGLE-3}

\end{axis}
\end{tikzpicture}
\caption{Зависимость acceptance rate от числа токенов драфтера $\gamma$}
\label{fig:acceptance_vs_gamma}
\end{figure}

% ------------------------------------------------------------
% ГРАФИК 3 — Ускорение vs γ (EAGLE-3 подробно)
% ------------------------------------------------------------
\begin{figure}[h]
\centering
\begin{tikzpicture}
\begin{axis}[
    width=13cm, height=7cm,
    xlabel={Число токенов драфтера $\gamma$},
    ylabel={Ускорение относительно target},
    legend pos=north west,
    legend style={font=\small},
    grid=major,
    grid style={dashed, gray!30},
    xmin=1, xmax=8,
    ymin=1.0, ymax=3.2,
    xtick={1,2,3,4,5,6,7,8},
    ytick={1.0,1.25,1.5,1.75,2.0,2.25,2.5,2.75,3.0},
    tick label style={font=\small},
    yticklabel style={/pgf/number format/fixed, /pgf/number format/precision=2},
]
\addplot[blue, thick, dashed, mark=square*] coordinates {
    (1,1.28)(2,1.46)(3,1.57)(4,1.66)(5,1.74)(6,1.79)(7,1.82)(8,1.82)
};
\addlegendentry{Arch4}

\addplot[purple, thick, mark=pentagon*] coordinates {
    (1,1.58)(2,1.88)(3,2.07)(4,2.27)(5,2.41)(6,2.51)(7,2.57)(8,2.57)
};
\addlegendentry{Arch1}

\addplot[red, very thick, mark=*] coordinates {
    (1,1.44)(2,1.84)(3,2.14)(4,2.42)(5,2.63)(6,2.72)(7,2.76)(8,2.74)
};
\addlegendentry{EAGLE-3}

% теоретический максимум β·γ+1 / 1 при β=0.761
\addplot[red, thin, dotted, domain=1:8, samples=50] {0.761*x + 1};
\addlegendentry{EAGLE-3 теор. макс. ($\beta \gamma + 1$)}

\end{axis}
\end{tikzpicture}
\caption{Зависимость ускорения инференса от числа токенов драфтера $\gamma$}
\label{fig:speedup_vs_gamma}
\end{figure}

% ------------------------------------------------------------
% ГРАФИК 4 — Столбчатый: итоговое сравнение архитектур
% ------------------------------------------------------------
\begin{figure}[h]
\centering
\begin{tikzpicture}
\begin{axis}[
    ybar,
    width=13cm, height=7cm,
    bar width=0.5cm,
    xlabel={Архитектура},
    ylabel={},
    legend pos=north west,
    legend style={font=\small},
    grid=major, grid style={dashed, gray!30},
    ymin=0, ymax=85,
    symbolic x coords={Arch4,Arch3,Arch2,Arch1,EAGLE-3},
    xtick=data,
    tick label style={font=\small},
    nodes near coords,
    nodes near coords style={font=\tiny},
    every node near coord/.append style={rotate=90, anchor=west},
]
\addplot[fill=blue!60] coordinates {
    (Arch4,43.7)(Arch3,52.4)(Arch2,61.8)(Arch1,68.3)(EAGLE-3,76.1)
};
\addlegendentry{Acceptance rate, \%}

\addplot[fill=red!60] coordinates {
    (Arch4,17.4)(Arch3,19.8)(Arch2,22.3)(Arch1,24.1)(EAGLE-3,26.3)
};
\addlegendentry{Ускорение $\times 10$}

\end{axis}
\end{tikzpicture}
\caption{Итоговое сравнение архитектур при $\gamma = 5$}
\label{fig:comparison_bar}
\end{figure}

% ------------------------------------------------------------
% ГРАФИК 5 — Feature loss и KL loss по шагам (EAGLE-3)
% ------------------------------------------------------------
\begin{figure}[h]
\centering
\begin{tikzpicture}
\begin{axis}[
    width=13cm, height=6.5cm,
    xlabel={Шаг обучения},
    ylabel={Loss},
    legend pos=north east,
    legend style={font=\small},
    grid=major,
    grid style={dashed, gray!30},
    xmin=0, xmax=4500,
    ymin=0, ymax=2.5,
    tick label style={font=\small},
]
\addplot[red, very thick] coordinates {
    (0,2.24)(200,1.58)(500,1.22)(1000,0.96)(1500,0.83)(2000,0.76)(2500,0.73)(3000,0.71)(3500,0.71)(4000,0.70)(4500,0.70)
};
\addlegendentry{KL distill loss}

\addplot[blue, thick, dashed] coordinates {
    (0,0.0)(200,0.0)(500,0.0)(1000,0.0)(1500,0.0)(2000,0.0)(2500,0.0)
};  % placeholder hidden
\addplot[blue, thick] coordinates {
    (0,0.58)(200,0.42)(500,0.31)(1000,0.21)(1500,0.14)(2000,0.08)(2500,0.06)(3000,0.04)(3500,0.03)(4000,0.03)(4500,0.03)
};
\addlegendentry{Feature alignment MSE}

\addplot[green!60!black, thick, dashed] coordinates {
    (0,2.41)(200,2.08)(500,1.87)(1000,1.63)(1500,1.49)(2000,1.41)(2500,1.37)(3000,1.35)(3500,1.34)(4000,1.33)(4500,1.33)
};
\addlegendentry{Task CE loss}

\end{axis}
\end{tikzpicture}
\caption{Динамика компонент функции потерь EAGLE-3 в процессе обучения}
\label{fig:eagle3_losses}
\end{figure}


\end{document}
